% Options for packages loaded elsewhere
\PassOptionsToPackage{unicode}{hyperref}
\PassOptionsToPackage{hyphens}{url}
\PassOptionsToPackage{dvipsnames,svgnames,x11names}{xcolor}
%
\documentclass[
]{article}
\usepackage{amsmath,amssymb}
\usepackage{iftex}
\ifPDFTeX
  \usepackage[T1]{fontenc}
  \usepackage[utf8]{inputenc}
  \usepackage{textcomp} % provide euro and other symbols
\else % if luatex or xetex
  \usepackage{unicode-math} % this also loads fontspec
  \defaultfontfeatures{Scale=MatchLowercase}
  \defaultfontfeatures[\rmfamily]{Ligatures=TeX,Scale=1}
\fi
\usepackage{lmodern}
\ifPDFTeX\else
  % xetex/luatex font selection
\fi
% Use upquote if available, for straight quotes in verbatim environments
\IfFileExists{upquote.sty}{\usepackage{upquote}}{}
\IfFileExists{microtype.sty}{% use microtype if available
  \usepackage[]{microtype}
  \UseMicrotypeSet[protrusion]{basicmath} % disable protrusion for tt fonts
}{}
\makeatletter
\@ifundefined{KOMAClassName}{% if non-KOMA class
  \IfFileExists{parskip.sty}{%
    \usepackage{parskip}
  }{% else
    \setlength{\parindent}{0pt}
    \setlength{\parskip}{6pt plus 2pt minus 1pt}}
}{% if KOMA class
  \KOMAoptions{parskip=half}}
\makeatother
\usepackage{xcolor}
\usepackage[margin=1in]{geometry}
\usepackage{color}
\usepackage{fancyvrb}
\newcommand{\VerbBar}{|}
\newcommand{\VERB}{\Verb[commandchars=\\\{\}]}
\DefineVerbatimEnvironment{Highlighting}{Verbatim}{commandchars=\\\{\}}
% Add ',fontsize=\small' for more characters per line
\usepackage{framed}
\definecolor{shadecolor}{RGB}{248,248,248}
\newenvironment{Shaded}{\begin{snugshade}}{\end{snugshade}}
\newcommand{\AlertTok}[1]{\textcolor[rgb]{0.94,0.16,0.16}{#1}}
\newcommand{\AnnotationTok}[1]{\textcolor[rgb]{0.56,0.35,0.01}{\textbf{\textit{#1}}}}
\newcommand{\AttributeTok}[1]{\textcolor[rgb]{0.13,0.29,0.53}{#1}}
\newcommand{\BaseNTok}[1]{\textcolor[rgb]{0.00,0.00,0.81}{#1}}
\newcommand{\BuiltInTok}[1]{#1}
\newcommand{\CharTok}[1]{\textcolor[rgb]{0.31,0.60,0.02}{#1}}
\newcommand{\CommentTok}[1]{\textcolor[rgb]{0.56,0.35,0.01}{\textit{#1}}}
\newcommand{\CommentVarTok}[1]{\textcolor[rgb]{0.56,0.35,0.01}{\textbf{\textit{#1}}}}
\newcommand{\ConstantTok}[1]{\textcolor[rgb]{0.56,0.35,0.01}{#1}}
\newcommand{\ControlFlowTok}[1]{\textcolor[rgb]{0.13,0.29,0.53}{\textbf{#1}}}
\newcommand{\DataTypeTok}[1]{\textcolor[rgb]{0.13,0.29,0.53}{#1}}
\newcommand{\DecValTok}[1]{\textcolor[rgb]{0.00,0.00,0.81}{#1}}
\newcommand{\DocumentationTok}[1]{\textcolor[rgb]{0.56,0.35,0.01}{\textbf{\textit{#1}}}}
\newcommand{\ErrorTok}[1]{\textcolor[rgb]{0.64,0.00,0.00}{\textbf{#1}}}
\newcommand{\ExtensionTok}[1]{#1}
\newcommand{\FloatTok}[1]{\textcolor[rgb]{0.00,0.00,0.81}{#1}}
\newcommand{\FunctionTok}[1]{\textcolor[rgb]{0.13,0.29,0.53}{\textbf{#1}}}
\newcommand{\ImportTok}[1]{#1}
\newcommand{\InformationTok}[1]{\textcolor[rgb]{0.56,0.35,0.01}{\textbf{\textit{#1}}}}
\newcommand{\KeywordTok}[1]{\textcolor[rgb]{0.13,0.29,0.53}{\textbf{#1}}}
\newcommand{\NormalTok}[1]{#1}
\newcommand{\OperatorTok}[1]{\textcolor[rgb]{0.81,0.36,0.00}{\textbf{#1}}}
\newcommand{\OtherTok}[1]{\textcolor[rgb]{0.56,0.35,0.01}{#1}}
\newcommand{\PreprocessorTok}[1]{\textcolor[rgb]{0.56,0.35,0.01}{\textit{#1}}}
\newcommand{\RegionMarkerTok}[1]{#1}
\newcommand{\SpecialCharTok}[1]{\textcolor[rgb]{0.81,0.36,0.00}{\textbf{#1}}}
\newcommand{\SpecialStringTok}[1]{\textcolor[rgb]{0.31,0.60,0.02}{#1}}
\newcommand{\StringTok}[1]{\textcolor[rgb]{0.31,0.60,0.02}{#1}}
\newcommand{\VariableTok}[1]{\textcolor[rgb]{0.00,0.00,0.00}{#1}}
\newcommand{\VerbatimStringTok}[1]{\textcolor[rgb]{0.31,0.60,0.02}{#1}}
\newcommand{\WarningTok}[1]{\textcolor[rgb]{0.56,0.35,0.01}{\textbf{\textit{#1}}}}
\usepackage{graphicx}
\makeatletter
\def\maxwidth{\ifdim\Gin@nat@width>\linewidth\linewidth\else\Gin@nat@width\fi}
\def\maxheight{\ifdim\Gin@nat@height>\textheight\textheight\else\Gin@nat@height\fi}
\makeatother
% Scale images if necessary, so that they will not overflow the page
% margins by default, and it is still possible to overwrite the defaults
% using explicit options in \includegraphics[width, height, ...]{}
\setkeys{Gin}{width=\maxwidth,height=\maxheight,keepaspectratio}
% Set default figure placement to htbp
\makeatletter
\def\fps@figure{htbp}
\makeatother
\setlength{\emergencystretch}{3em} % prevent overfull lines
\providecommand{\tightlist}{%
  \setlength{\itemsep}{0pt}\setlength{\parskip}{0pt}}
\setcounter{secnumdepth}{-\maxdimen} % remove section numbering
\ifLuaTeX
  \usepackage{selnolig}  % disable illegal ligatures
\fi
\usepackage{bookmark}
\IfFileExists{xurl.sty}{\usepackage{xurl}}{} % add URL line breaks if available
\urlstyle{same}
\hypersetup{
  pdftitle={Lab 02},
  colorlinks=true,
  linkcolor={Maroon},
  filecolor={Maroon},
  citecolor={Blue},
  urlcolor={Blue},
  pdfcreator={LaTeX via pandoc}}

\title{Lab 02}
\author{}
\date{\vspace{-2.5em}}

\begin{document}
\maketitle

\section{Lab 02: R as a GIS}\label{lab-02-r-as-a-gis}

\subsubsection{Read the instructions COMPLETELY before starting the
lab}\label{read-the-instructions-completely-before-starting-the-lab}

This lab builds on many of the discussions and exercises from class,
including previous labs.

\subsubsection{Formatting your
submission}\label{formatting-your-submission}

This lab must be placed into a public repository on GitHub
(www.github.com). Before the due date, submit \textbf{on Canvas} a link
to the repository. I will then download your repositories and run your
code. The code must be contained in either a .R script or a .Rmd
markdown document. As I need to run your code, any data you use in the
lab must be referenced using \textbf{relative path names}. Finally,
answers to questions I pose in this document must also be in the
repository at the time you submit your link to Canvas. They can be in a
separate text file, or if you decide to use an RMarkdown document, you
can answer them directly in the doc.

\subsection{Data}\label{data}

The data for this lab can be found in the \texttt{./data/CBW/} directory
within the course GitHub repository.

Spatial datasets:

\begin{enumerate}
\def\labelenumi{\arabic{enumi}.}
\item
  Streams\_Opened\_by\_Dam\_Removal\_2012\_2017.shp
\item
  Dam\_or\_Other\_Blockage\_Removed\_2012\_2017.shp
\item
  County\_Boundaries.shp
\end{enumerate}

Non-spatial datasets

\begin{enumerate}
\def\labelenumi{\arabic{enumi}.}
\tightlist
\item
  BMPreport2016\_landbmps.csv
\end{enumerate}

\subsection{Working with tabular data}\label{working-with-tabular-data}

\begin{Shaded}
\begin{Highlighting}[]
\CommentTok{\# setup}
\FunctionTok{library}\NormalTok{(tidyverse)}
\end{Highlighting}
\end{Shaded}

\begin{verbatim}
## Warning: package 'tidyverse' was built under R version 4.4.2
\end{verbatim}

\begin{verbatim}
## Warning: package 'ggplot2' was built under R version 4.4.3
\end{verbatim}

\begin{verbatim}
## Warning: package 'readr' was built under R version 4.4.2
\end{verbatim}

\begin{verbatim}
## Warning: package 'purrr' was built under R version 4.4.3
\end{verbatim}

\begin{verbatim}
## Warning: package 'dplyr' was built under R version 4.4.3
\end{verbatim}

\begin{verbatim}
## Warning: package 'forcats' was built under R version 4.4.2
\end{verbatim}

\begin{verbatim}
## Warning: package 'lubridate' was built under R version 4.4.2
\end{verbatim}

\begin{verbatim}
## -- Attaching core tidyverse packages ------------------------ tidyverse 2.0.0 --
## v dplyr     1.1.4     v readr     2.1.5
## v forcats   1.0.0     v stringr   1.5.1
## v ggplot2   4.0.0     v tibble    3.2.1
## v lubridate 1.9.4     v tidyr     1.3.1
## v purrr     1.1.0     
## -- Conflicts ------------------------------------------ tidyverse_conflicts() --
## x dplyr::filter() masks stats::filter()
## x dplyr::lag()    masks stats::lag()
## i Use the conflicted package (<http://conflicted.r-lib.org/>) to force all conflicts to become errors
\end{verbatim}

\begin{Shaded}
\begin{Highlighting}[]
\FunctionTok{library}\NormalTok{(sf)}
\end{Highlighting}
\end{Shaded}

\begin{verbatim}
## Warning: package 'sf' was built under R version 4.4.2
\end{verbatim}

\begin{verbatim}
## Linking to GEOS 3.12.2, GDAL 3.9.3, PROJ 9.4.1; sf_use_s2() is TRUE
\end{verbatim}

\begin{Shaded}
\begin{Highlighting}[]
\FunctionTok{library}\NormalTok{(tmap)}
\end{Highlighting}
\end{Shaded}

\begin{verbatim}
## Warning: package 'tmap' was built under R version 4.4.3
\end{verbatim}

A ``join'' is a method to join multiple tables together using a matching
``key'' found in both datasets. For example:

\begin{Shaded}
\begin{Highlighting}[]
\CommentTok{\# table of names}
\NormalTok{t.names }\OtherTok{\textless{}{-}} \FunctionTok{tibble}\NormalTok{(}\AttributeTok{key =} \FunctionTok{c}\NormalTok{(}\DecValTok{1}\NormalTok{, }\DecValTok{2}\NormalTok{, }\DecValTok{3}\NormalTok{), }
             \AttributeTok{name =} \FunctionTok{c}\NormalTok{(}\StringTok{"Huey"}\NormalTok{, }\StringTok{"Dewey"}\NormalTok{, }\StringTok{"Louis"}\NormalTok{))}

\CommentTok{\# table of scores}
\NormalTok{t.scores }\OtherTok{\textless{}{-}} \FunctionTok{tibble}\NormalTok{(}\AttributeTok{name =} \FunctionTok{c}\NormalTok{(}\StringTok{"Louis"}\NormalTok{, }\StringTok{"Huey"}\NormalTok{, }\StringTok{"Dewey"}\NormalTok{),}
                   \AttributeTok{grade =} \FunctionTok{c}\NormalTok{(}\DecValTok{99}\NormalTok{, }\DecValTok{45}\NormalTok{, }\DecValTok{33}\NormalTok{))}

\CommentTok{\# join them together using a common variable, in this case the "name" variable}
\NormalTok{t.joined }\OtherTok{\textless{}{-}} \FunctionTok{left\_join}\NormalTok{(t.names, t.scores, }\AttributeTok{by =} \StringTok{"name"}\NormalTok{)}
\NormalTok{t.joined}
\end{Highlighting}
\end{Shaded}

\begin{verbatim}
## # A tibble: 3 x 3
##     key name  grade
##   <dbl> <chr> <dbl>
## 1     1 Huey     45
## 2     2 Dewey    33
## 3     3 Louis    99
\end{verbatim}

A ``left join'' finds starts with the table on the ``left'' and then
finds matches in the table on the ``right''. See the documentation using
\texttt{?left\_join} for more details and for other types of joins.
Sometimes the attributes you're using to join the tables won't have the
same name, in which case the you need to tell the join function which
variables to join on. Note the difference in the \texttt{by} argument in
the below code block compared to the one above.

\begin{Shaded}
\begin{Highlighting}[]
\NormalTok{t.wonkyNames }\OtherTok{\textless{}{-}} \FunctionTok{tibble}\NormalTok{(}\AttributeTok{nombre =} \FunctionTok{c}\NormalTok{(}\StringTok{"Dewey"}\NormalTok{, }\StringTok{"Louis"}\NormalTok{, }\StringTok{"Huey"}\NormalTok{),}
                       \AttributeTok{x =} \FunctionTok{rep}\NormalTok{(}\DecValTok{999}\NormalTok{),}
                       \AttributeTok{favoriteFood =} \FunctionTok{c}\NormalTok{(}\StringTok{"banana"}\NormalTok{, }\StringTok{"apple"}\NormalTok{, }\StringTok{"carrot"}\NormalTok{))}

\NormalTok{t.joined2 }\OtherTok{\textless{}{-}} \FunctionTok{left\_join}\NormalTok{(t.names, t.wonkyNames, }\AttributeTok{by =} \FunctionTok{c}\NormalTok{(}\StringTok{"name"} \OtherTok{=} \StringTok{"nombre"}\NormalTok{))}
\NormalTok{t.joined2}
\end{Highlighting}
\end{Shaded}

\begin{verbatim}
## # A tibble: 3 x 4
##     key name      x favoriteFood
##   <dbl> <chr> <dbl> <chr>       
## 1     1 Huey    999 carrot      
## 2     2 Dewey   999 banana      
## 3     3 Louis   999 apple
\end{verbatim}

\subsection{Let's take a look at some tabular
data}\label{lets-take-a-look-at-some-tabular-data}

This dataset includes a list of best management practices (``BMPs'') to
reduce nutrient and sediment pollution in the Chesapeake Bay Watershed.

\begin{Shaded}
\begin{Highlighting}[]
\NormalTok{bmps }\OtherTok{\textless{}{-}} \FunctionTok{read\_csv}\NormalTok{(}\StringTok{"../Lab\_2/CBW/BMPreport2016\_landbmps.csv"}\NormalTok{)}
\end{Highlighting}
\end{Shaded}

\begin{verbatim}
## Rows: 69601 Columns: 18
## -- Column specification --------------------------------------------------------
## Delimiter: ","
## chr (11): StateAbbreviation, GeographyName, Geography, Agency, BMPShortName,...
## dbl  (7): AmountSubmitted, AmountBackedOut, AmountNotBackedOut, AmountCredit...
## 
## i Use `spec()` to retrieve the full column specification for this data.
## i Specify the column types or set `show_col_types = FALSE` to quiet this message.
\end{verbatim}

\begin{Shaded}
\begin{Highlighting}[]
\FunctionTok{glimpse}\NormalTok{(bmps)}
\end{Highlighting}
\end{Shaded}

\begin{verbatim}
## Rows: 69,601
## Columns: 18
## $ StateAbbreviation   <chr> "DC", "DC", "DC", "DC", "DC", "DC", "DC", "DC", "D~
## $ GeographyName       <chr> "11001(cbwsonly)", "11001(cbwsonly)", "11001(cbwso~
## $ Geography           <chr> "Washington, DC (CBWS Portion Only)", "Washington,~
## $ Agency              <chr> "Department of Defense", "Department of Defense", ~
## $ BMPShortName        <chr> "wetpondwetland", "wetpondwetland", "wetpondwetlan~
## $ BMP                 <chr> "Wet Ponds and Wetlands", "Wet Ponds and Wetlands"~
## $ BMPType             <chr> "Efficiency", "Efficiency", "Efficiency", "Efficie~
## $ Unit                <chr> "Acres Treated", "Acres Treated", "Acres Treated",~
## $ Sector              <chr> "Developed", "Developed", "Developed", "Developed"~
## $ FromLoadSource      <chr> "Non-Regulated Roads", "Non-Regulated Buildings an~
## $ ToLoadSource        <chr> "Non-Regulated Roads", "Non-Regulated Buildings an~
## $ AmountSubmitted     <dbl> 8.709123261, 47.984171150, 1.358309392, 6.76626660~
## $ AmountBackedOut     <dbl> 0, 0, 0, 0, 0, 0, 0, 0, 0, 0, 0, 0, 0, 0, 0, 0, 0,~
## $ AmountNotBackedOut  <dbl> 8.709123261, 47.984171150, 1.358309392, 6.76626660~
## $ AmountCredited      <dbl> 8.709123261, 47.984171150, 1.358309392, 6.76626660~
## $ Excess              <dbl> 0, 0, 0, 0, 0, 0, 0, 0, 0, 0, 0, 0, 0, 0, 0, 0, 0,~
## $ TotalAmountCredited <dbl> 8.709123261, 47.984171150, 1.358309392, 6.76626660~
## $ Cost                <dbl> 11462.077120, 63151.967650, 1787.670991, 8905.0834~
\end{verbatim}

Look at the attribute ``GeographyName'' - it's a character attribute
that contains the counties' FIPS code, but also some ancillary
explanatory data we need to get rid of. There are multiple ways of doing
so, including some (very) fancy automated methods that detect patterns
of numbers and characters. We're going to take a simpler approach and
assume that all FIPS codes are only 5 characters long.

\begin{Shaded}
\begin{Highlighting}[]
\CommentTok{\# edit the bmps variable in place, which isn\textquotesingle{}t always best practices}
\NormalTok{bmps }\OtherTok{\textless{}{-}}\NormalTok{ bmps }\SpecialCharTok{\%\textgreater{}\%} \FunctionTok{mutate}\NormalTok{(., }\AttributeTok{FIPS.trimmed =}\NormalTok{ stringr}\SpecialCharTok{::}\FunctionTok{str\_sub}\NormalTok{(GeographyName, }\DecValTok{1}\NormalTok{, }\DecValTok{5}\NormalTok{))}
\end{Highlighting}
\end{Shaded}

This preparing raw data for further analysis is commonly called
``munging''. While the term isn't the important part, these kinds of
preparatory steps can be useful when you're trying to create ``keys'' by
which to join tables or just clean your tables in general.

Let's recall how to do some simple tasks

\begin{Shaded}
\begin{Highlighting}[]
\CommentTok{\# Let\textquotesingle{}s calculate the total cost by BMP and then plot it}
\NormalTok{bmps }\SpecialCharTok{\%\textgreater{}\%} \FunctionTok{group\_by}\NormalTok{(BMPType) }\SpecialCharTok{\%\textgreater{}\%} \FunctionTok{summarise}\NormalTok{(}\AttributeTok{totalCost =} \FunctionTok{sum}\NormalTok{(Cost)) }\SpecialCharTok{\%\textgreater{}\%}
  \FunctionTok{ggplot}\NormalTok{(., }\FunctionTok{aes}\NormalTok{(}\AttributeTok{x =}\NormalTok{ BMPType, }\AttributeTok{y =}\NormalTok{ totalCost)) }\SpecialCharTok{+}
  \FunctionTok{geom\_bar}\NormalTok{(}\AttributeTok{stat =} \StringTok{"identity"}\NormalTok{) }\SpecialCharTok{+}
  \FunctionTok{theme\_minimal}\NormalTok{()}
\end{Highlighting}
\end{Shaded}

\begin{verbatim}
## Warning: Removed 2 rows containing missing values or values outside the scale range
## (`geom_bar()`).
\end{verbatim}

\includegraphics{lab02_files/figure-latex/review-1.pdf}

\begin{Shaded}
\begin{Highlighting}[]
\CommentTok{\# Doesn\textquotesingle{}t really work. This is because there are missing data in the cost attribute. }
\CommentTok{\# Let\textquotesingle{}s look at it (and this is why we do exploratory data analysis first)}

\FunctionTok{summary}\NormalTok{(bmps}\SpecialCharTok{$}\NormalTok{Cost)}
\end{Highlighting}
\end{Shaded}

\begin{verbatim}
##     Min.  1st Qu.   Median     Mean  3rd Qu.     Max.     NA's 
##        0        0      142    37408     3583 39731974    14286
\end{verbatim}

\begin{Shaded}
\begin{Highlighting}[]
\CommentTok{\# Yup, lots of "NA\textquotesingle{}s"... We can drop them in our analysis. }
\CommentTok{\# Look carefully at the code in the \textquotesingle{}sum\textquotesingle{} function. What is the na.rm = T parameter doing?}


\NormalTok{bmps }\SpecialCharTok{\%\textgreater{}\%} \FunctionTok{group\_by}\NormalTok{(BMPType) }\SpecialCharTok{\%\textgreater{}\%} \FunctionTok{summarise}\NormalTok{(}\AttributeTok{totalCost =} \FunctionTok{sum}\NormalTok{(Cost, }\AttributeTok{na.rm =}\NormalTok{ T)) }\SpecialCharTok{\%\textgreater{}\%}
  \FunctionTok{ggplot}\NormalTok{(., }\FunctionTok{aes}\NormalTok{(}\AttributeTok{x =}\NormalTok{ BMPType, }\AttributeTok{y =}\NormalTok{ totalCost)) }\SpecialCharTok{+}
  \FunctionTok{geom\_bar}\NormalTok{(}\AttributeTok{stat =} \StringTok{"identity"}\NormalTok{) }\SpecialCharTok{+}
  \FunctionTok{theme\_minimal}\NormalTok{()}
\end{Highlighting}
\end{Shaded}

\includegraphics{lab02_files/figure-latex/review-2.pdf}

We can also group by multiple variables at the same time. For example:

\begin{Shaded}
\begin{Highlighting}[]
\CommentTok{\# group by state and sector, sum total cost}
\NormalTok{twofactors }\OtherTok{\textless{}{-}}\NormalTok{ bmps }\SpecialCharTok{\%\textgreater{}\%} \FunctionTok{group\_by}\NormalTok{(StateAbbreviation, Sector) }\SpecialCharTok{\%\textgreater{}\%} 
  \FunctionTok{summarise}\NormalTok{(}\AttributeTok{totalCost =} \FunctionTok{sum}\NormalTok{(Cost))}
\end{Highlighting}
\end{Shaded}

\begin{verbatim}
## `summarise()` has grouped output by 'StateAbbreviation'. You can override using
## the `.groups` argument.
\end{verbatim}

\begin{Shaded}
\begin{Highlighting}[]
\NormalTok{twofactors}
\end{Highlighting}
\end{Shaded}

\begin{verbatim}
## # A tibble: 25 x 3
## # Groups:   StateAbbreviation [7]
##    StateAbbreviation Sector      totalCost
##    <chr>             <chr>           <dbl>
##  1 DC                Developed         NA 
##  2 DC                Natural           NA 
##  3 DE                Agriculture       NA 
##  4 DE                Developed         NA 
##  5 DE                Natural           NA 
##  6 DE                Septic       2064477.
##  7 MD                Agriculture       NA 
##  8 MD                Developed         NA 
##  9 MD                Natural           NA 
## 10 MD                Septic      50470281.
## # i 15 more rows
\end{verbatim}

For our last bit of review, let's make a few box plots:

\begin{Shaded}
\begin{Highlighting}[]
\CommentTok{\# A simple one}
\NormalTok{bmps }\SpecialCharTok{\%\textgreater{}\%} \FunctionTok{ggplot}\NormalTok{(., }\FunctionTok{aes}\NormalTok{(}\AttributeTok{x =}\NormalTok{ StateAbbreviation, }\AttributeTok{y =}\NormalTok{ AmountCredited)) }\SpecialCharTok{+}
  \FunctionTok{geom\_boxplot}\NormalTok{(}\FunctionTok{aes}\NormalTok{(}\AttributeTok{fill =}\NormalTok{ StateAbbreviation))}
\end{Highlighting}
\end{Shaded}

\includegraphics{lab02_files/figure-latex/review3-1.pdf}

\begin{Shaded}
\begin{Highlighting}[]
\CommentTok{\# Very heavily skewed, so just for the sake of visualization, }
\CommentTok{\# let\textquotesingle{}s subset the data (dramatically)}
\NormalTok{bmps }\SpecialCharTok{\%\textgreater{}\%} 
\NormalTok{  dplyr}\SpecialCharTok{::}\FunctionTok{filter}\NormalTok{(., AmountCredited }\SpecialCharTok{\textgreater{}} \DecValTok{1} \SpecialCharTok{\&}\NormalTok{ AmountCredited }\SpecialCharTok{\textless{}} \DecValTok{100}\NormalTok{) }\SpecialCharTok{\%\textgreater{}\%} 
  \FunctionTok{ggplot}\NormalTok{(., }\FunctionTok{aes}\NormalTok{(}\AttributeTok{x =}\NormalTok{ StateAbbreviation, }\AttributeTok{y =}\NormalTok{ AmountCredited)) }\SpecialCharTok{+}
  \FunctionTok{geom\_boxplot}\NormalTok{(}\FunctionTok{aes}\NormalTok{(}\AttributeTok{fill =}\NormalTok{ StateAbbreviation))}
\end{Highlighting}
\end{Shaded}

\includegraphics{lab02_files/figure-latex/review3-2.pdf}

\begin{Shaded}
\begin{Highlighting}[]
\CommentTok{\# We can also plot multiple dimensions in our plot using }
\CommentTok{\# the \textasciigrave{}facet\textasciigrave{} family of commands in ggplot}

\NormalTok{bmps }\SpecialCharTok{\%\textgreater{}\%} 
\NormalTok{  dplyr}\SpecialCharTok{::}\FunctionTok{filter}\NormalTok{(., AmountCredited }\SpecialCharTok{\textgreater{}} \DecValTok{1} \SpecialCharTok{\&}\NormalTok{ AmountCredited }\SpecialCharTok{\textless{}} \DecValTok{100}\NormalTok{) }\SpecialCharTok{\%\textgreater{}\%} 
  \FunctionTok{ggplot}\NormalTok{(., }\FunctionTok{aes}\NormalTok{(}\AttributeTok{x =}\NormalTok{ StateAbbreviation, }\AttributeTok{y =}\NormalTok{ AmountCredited)) }\SpecialCharTok{+}
  \FunctionTok{geom\_boxplot}\NormalTok{(}\FunctionTok{aes}\NormalTok{(}\AttributeTok{fill =}\NormalTok{ StateAbbreviation)) }\SpecialCharTok{+}
  \FunctionTok{facet\_grid}\NormalTok{(Sector}\SpecialCharTok{\textasciitilde{}}\NormalTok{.)}
\end{Highlighting}
\end{Shaded}

\includegraphics{lab02_files/figure-latex/review3-3.pdf}

\begin{Shaded}
\begin{Highlighting}[]
\CommentTok{\# finally, an alternative to removing the outliers would be to log{-}transform the y{-}axis}
\CommentTok{\# But be careful, this removes a lot of the data because the log of 0 is undefined}
\NormalTok{bmps }\SpecialCharTok{\%\textgreater{}\%} \FunctionTok{ggplot}\NormalTok{(., }\FunctionTok{aes}\NormalTok{(}\AttributeTok{x =}\NormalTok{ StateAbbreviation, }\AttributeTok{y =}\NormalTok{ AmountCredited)) }\SpecialCharTok{+}
  \FunctionTok{geom\_boxplot}\NormalTok{(}\FunctionTok{aes}\NormalTok{(}\AttributeTok{fill =}\NormalTok{ StateAbbreviation)) }\SpecialCharTok{+}
  \FunctionTok{scale\_y\_log10}\NormalTok{() }\SpecialCharTok{+}
  \FunctionTok{labs}\NormalTok{(}\AttributeTok{y =} \StringTok{"log10(AmountCredited)"}\NormalTok{)}
\end{Highlighting}
\end{Shaded}

\begin{verbatim}
## Warning in scale_y_log10(): log-10 transformation introduced infinite values.
\end{verbatim}

\begin{verbatim}
## Warning: Removed 15585 rows containing non-finite outside the scale range
## (`stat_boxplot()`).
\end{verbatim}

\includegraphics{lab02_files/figure-latex/review3-4.pdf}

The last new item uses the \texttt{\%in\%} command. It's a way to
quickly figure out which elements are inside of another. In that sense,
it's similar to a spatial intersection, but for other types of data.

\begin{Shaded}
\begin{Highlighting}[]
\CommentTok{\# setup a demonstration vector}
\NormalTok{x }\OtherTok{\textless{}{-}} \FunctionTok{c}\NormalTok{(}\DecValTok{1}\NormalTok{, }\DecValTok{2}\NormalTok{, }\DecValTok{3}\NormalTok{, }\DecValTok{4}\NormalTok{, }\DecValTok{5}\NormalTok{)}

\CommentTok{\# is 7 in our vector?}
\DecValTok{7} \SpecialCharTok{\%in\%}\NormalTok{ x }\CommentTok{\# should evaluate False}
\end{Highlighting}
\end{Shaded}

\begin{verbatim}
## [1] FALSE
\end{verbatim}

\begin{Shaded}
\begin{Highlighting}[]
\DecValTok{2} \SpecialCharTok{\%in\%}\NormalTok{ x }\CommentTok{\# should evaluate True}
\end{Highlighting}
\end{Shaded}

\begin{verbatim}
## [1] TRUE
\end{verbatim}

\begin{Shaded}
\begin{Highlighting}[]
\CommentTok{\# can also do it with vectors}
\FunctionTok{c}\NormalTok{(}\DecValTok{4}\NormalTok{, }\DecValTok{99}\NormalTok{, }\DecValTok{1}\NormalTok{) }\SpecialCharTok{\%in\%}\NormalTok{ x}
\end{Highlighting}
\end{Shaded}

\begin{verbatim}
## [1]  TRUE FALSE  TRUE
\end{verbatim}

Lastly, let's recall using tmap on our data. Also remember you can use
\texttt{sf::st\_make\_valid} to fix offending geometry

\begin{Shaded}
\begin{Highlighting}[]
\NormalTok{counties }\OtherTok{\textless{}{-}}\NormalTok{ sf}\SpecialCharTok{::}\FunctionTok{read\_sf}\NormalTok{(}\StringTok{"../Lab\_2/CBW/County\_Boundaries.shp"}\NormalTok{)}
\NormalTok{counties }\SpecialCharTok{\%\textgreater{}\%}\NormalTok{ sf}\SpecialCharTok{::}\FunctionTok{st\_is\_valid}\NormalTok{()}
\end{Highlighting}
\end{Shaded}

\begin{verbatim}
##   [1]  TRUE  TRUE  TRUE  TRUE  TRUE  TRUE  TRUE  TRUE  TRUE  TRUE  TRUE  TRUE
##  [13]  TRUE  TRUE  TRUE  TRUE  TRUE  TRUE  TRUE  TRUE  TRUE  TRUE  TRUE  TRUE
##  [25]  TRUE  TRUE  TRUE  TRUE  TRUE  TRUE  TRUE  TRUE  TRUE  TRUE  TRUE  TRUE
##  [37]  TRUE  TRUE  TRUE  TRUE  TRUE  TRUE  TRUE  TRUE  TRUE  TRUE  TRUE  TRUE
##  [49]  TRUE  TRUE  TRUE  TRUE  TRUE  TRUE  TRUE  TRUE  TRUE  TRUE  TRUE  TRUE
##  [61]  TRUE  TRUE  TRUE  TRUE  TRUE  TRUE  TRUE  TRUE  TRUE  TRUE  TRUE  TRUE
##  [73]  TRUE  TRUE  TRUE  TRUE  TRUE  TRUE  TRUE  TRUE  TRUE  TRUE  TRUE  TRUE
##  [85]  TRUE  TRUE  TRUE  TRUE  TRUE  TRUE  TRUE  TRUE  TRUE  TRUE  TRUE  TRUE
##  [97]  TRUE  TRUE  TRUE  TRUE  TRUE  TRUE  TRUE  TRUE  TRUE  TRUE  TRUE  TRUE
## [109]  TRUE  TRUE  TRUE  TRUE  TRUE  TRUE  TRUE  TRUE  TRUE  TRUE  TRUE  TRUE
## [121]  TRUE  TRUE  TRUE  TRUE  TRUE  TRUE  TRUE  TRUE  TRUE  TRUE  TRUE  TRUE
## [133]  TRUE  TRUE  TRUE  TRUE  TRUE  TRUE  TRUE  TRUE  TRUE  TRUE  TRUE  TRUE
## [145]  TRUE  TRUE  TRUE  TRUE  TRUE  TRUE  TRUE  TRUE  TRUE  TRUE  TRUE  TRUE
## [157]  TRUE  TRUE  TRUE  TRUE  TRUE  TRUE  TRUE  TRUE  TRUE  TRUE  TRUE  TRUE
## [169]  TRUE  TRUE  TRUE  TRUE  TRUE  TRUE  TRUE  TRUE  TRUE  TRUE  TRUE  TRUE
## [181]  TRUE  TRUE  TRUE  TRUE  TRUE  TRUE  TRUE  TRUE  TRUE  TRUE  TRUE  TRUE
## [193]  TRUE  TRUE  TRUE  TRUE  TRUE  TRUE  TRUE  TRUE  TRUE  TRUE  TRUE  TRUE
## [205]  TRUE FALSE  TRUE
\end{verbatim}

\begin{Shaded}
\begin{Highlighting}[]
\NormalTok{counties }\OtherTok{\textless{}{-}}\NormalTok{ counties }\SpecialCharTok{\%\textgreater{}\%}\NormalTok{ sf}\SpecialCharTok{::}\FunctionTok{st\_make\_valid}\NormalTok{()}

\CommentTok{\# quick map of the data}
\FunctionTok{tm\_shape}\NormalTok{(counties) }\SpecialCharTok{+} \FunctionTok{tm\_polygons}\NormalTok{(}\AttributeTok{fill =} \StringTok{"ALAND10"}\NormalTok{)}
\end{Highlighting}
\end{Shaded}

\includegraphics{lab02_files/figure-latex/remembering tmap-1.pdf}

\newpage

\subsection{Your tasks}\label{your-tasks}

Using the following data\ldots{}

\begin{Shaded}
\begin{Highlighting}[]
\CommentTok{\# spatial}
\CommentTok{\# Note, the st\_make\_valid() calls are NOT necessary in all cases}
\CommentTok{\# But we know it\textquotesingle{}s needed for the counties, and doesn\textquotesingle{}t really hurt anything}
\CommentTok{\# for the others}
\NormalTok{counties }\OtherTok{\textless{}{-}}\NormalTok{ sf}\SpecialCharTok{::}\FunctionTok{read\_sf}\NormalTok{(}\StringTok{"../Lab\_2/CBW/County\_Boundaries.shp"}\NormalTok{) }\SpecialCharTok{\%\textgreater{}\%} 
\NormalTok{  sf}\SpecialCharTok{::}\FunctionTok{st\_make\_valid}\NormalTok{()}
\NormalTok{dams }\OtherTok{\textless{}{-}}\NormalTok{ sf}\SpecialCharTok{::}\FunctionTok{read\_sf}\NormalTok{(}\StringTok{"../Lab\_2/CBW/Dam\_or\_Other\_Blockage\_Removed\_2012\_2017.shp"}\NormalTok{) }\SpecialCharTok{\%\textgreater{}\%} 
\NormalTok{  sf}\SpecialCharTok{::}\FunctionTok{st\_make\_valid}\NormalTok{()}
\NormalTok{streams }\OtherTok{\textless{}{-}}\NormalTok{ sf}\SpecialCharTok{::}\FunctionTok{read\_sf}\NormalTok{(}\StringTok{"../Lab\_2/CBW/Streams\_Opened\_by\_Dam\_Removal\_2012\_2017.shp"}\NormalTok{) }\SpecialCharTok{\%\textgreater{}\%} 
\NormalTok{  sf}\SpecialCharTok{::}\FunctionTok{st\_make\_valid}\NormalTok{()}

\CommentTok{\# aspatial}
\NormalTok{bmps }\OtherTok{\textless{}{-}} \FunctionTok{read\_csv}\NormalTok{(}\StringTok{"../Lab\_2/CBW/BMPreport2016\_landbmps.csv"}\NormalTok{)}
\end{Highlighting}
\end{Shaded}

\begin{verbatim}
## Rows: 69601 Columns: 18
## -- Column specification --------------------------------------------------------
## Delimiter: ","
## chr (11): StateAbbreviation, GeographyName, Geography, Agency, BMPShortName,...
## dbl  (7): AmountSubmitted, AmountBackedOut, AmountNotBackedOut, AmountCredit...
## 
## i Use `spec()` to retrieve the full column specification for this data.
## i Specify the column types or set `show_col_types = FALSE` to quiet this message.
\end{verbatim}

\ldots complete the following tasks.

Complete each task COMPLETELY USING R CODE. YOU MUST SHOW YOUR WORK FOR
EACH ANSWER. Label your variables sensibly and use comments such that I
can find your answers and your work. The following tasks draw upon
lecture, your assigned readings, and the examples shown above. As
always, there are multiple ways of completing each task. Remember, it's
always a good idea to perform some exploratory data analysis on your own
prior to starting work.

\subsubsection{Task 1: Aspatial
operations}\label{task-1-aspatial-operations}

1.1 Calculate summary statistics for the Cost of BMPs for each State
(including DC)

\begin{Shaded}
\begin{Highlighting}[]
\CommentTok{\#summary stats including DC}
\FunctionTok{by}\NormalTok{(bmps}\SpecialCharTok{$}\NormalTok{Cost, bmps}\SpecialCharTok{$}\NormalTok{StateAbbreviation, summary)}
\end{Highlighting}
\end{Shaded}

\begin{verbatim}
## bmps$StateAbbreviation: DC
##      Min.   1st Qu.    Median      Mean   3rd Qu.      Max.      NA's 
##       0.1      26.7     628.6   22192.8    5000.7 1891063.7         8 
## ------------------------------------------------------------ 
## bmps$StateAbbreviation: DE
##    Min. 1st Qu.  Median    Mean 3rd Qu.    Max.    NA's 
##       0      16     906   30682   11574 5041838     238 
## ------------------------------------------------------------ 
## bmps$StateAbbreviation: MD
##     Min.  1st Qu.   Median     Mean  3rd Qu.     Max.     NA's 
##        0        2      876    53118    10229 12086857     2151 
## ------------------------------------------------------------ 
## bmps$StateAbbreviation: NY
##     Min.  1st Qu.   Median     Mean  3rd Qu.     Max.     NA's 
##      0.0    335.1   1954.4   8128.7   6978.7 622138.0      721 
## ------------------------------------------------------------ 
## bmps$StateAbbreviation: PA
##     Min.  1st Qu.   Median     Mean  3rd Qu.     Max.     NA's 
##        0        0      431    59904     5954 19968186     3565 
## ------------------------------------------------------------ 
## bmps$StateAbbreviation: VA
##     Min.  1st Qu.   Median     Mean  3rd Qu.     Max.     NA's 
##        0        0       15    25380     1240 39731974     7002 
## ------------------------------------------------------------ 
## bmps$StateAbbreviation: WV
##     Min.  1st Qu.   Median     Mean  3rd Qu.     Max.     NA's 
##        0        0       97    50014     1176 15796521      601
\end{verbatim}

1.2 Make a scatterplot of Cost vs.~TotalAmountCredited, ONLY FOR BMPs
with \texttt{Units} of type ``Acres''. You may need to apply a data
transformation to one or more axes if the data are heavily skewed.

\begin{Shaded}
\begin{Highlighting}[]
\CommentTok{\#select BMPs with units of Acres}
\NormalTok{Acres\_bmps }\OtherTok{\textless{}{-}}\NormalTok{ bmps }\SpecialCharTok{|\textgreater{}}
  \FunctionTok{filter}\NormalTok{(Unit}\SpecialCharTok{==}\StringTok{"Acres"}\NormalTok{)}

\CommentTok{\#make scatterplot}
\NormalTok{acre\_scatter }\OtherTok{\textless{}{-}}
  \FunctionTok{ggplot}\NormalTok{(Acres\_bmps, }\FunctionTok{aes}\NormalTok{(}\AttributeTok{x =}\NormalTok{ Cost, }\AttributeTok{y =}\NormalTok{ TotalAmountCredited)) }\SpecialCharTok{+}
  \FunctionTok{geom\_point}\NormalTok{()}

\NormalTok{acre\_scatter}
\end{Highlighting}
\end{Shaded}

\begin{verbatim}
## Warning: Removed 170 rows containing missing values or values outside the scale range
## (`geom_point()`).
\end{verbatim}

\includegraphics{lab02_files/figure-latex/unnamed-chunk-2-1.pdf}

\begin{Shaded}
\begin{Highlighting}[]
\CommentTok{\#need to transform this data{-} starting with y axis}
\NormalTok{y\_log\_scatter }\OtherTok{\textless{}{-}}
\NormalTok{  Acres\_bmps }\SpecialCharTok{\%\textgreater{}\%}
    \FunctionTok{ggplot}\NormalTok{(., }\FunctionTok{aes}\NormalTok{(}\AttributeTok{x =}\NormalTok{ Cost, }\AttributeTok{y =}\NormalTok{ TotalAmountCredited)) }\SpecialCharTok{+}
    \FunctionTok{geom\_point}\NormalTok{() }\SpecialCharTok{+}
    \FunctionTok{scale\_y\_log10}\NormalTok{() }\SpecialCharTok{+}
    \FunctionTok{labs}\NormalTok{(}\AttributeTok{y =} \StringTok{"log10(TotalAmountCredited)"}\NormalTok{)}

\NormalTok{y\_log\_scatter}
\end{Highlighting}
\end{Shaded}

\begin{verbatim}
## Warning in scale_y_log10(): log-10 transformation introduced infinite values.
## Removed 170 rows containing missing values or values outside the scale range
## (`geom_point()`).
\end{verbatim}

\includegraphics{lab02_files/figure-latex/unnamed-chunk-2-2.pdf}

\begin{Shaded}
\begin{Highlighting}[]
\CommentTok{\#trying x axis instead}
\NormalTok{x\_log\_scatter }\OtherTok{\textless{}{-}}
\NormalTok{  Acres\_bmps }\SpecialCharTok{\%\textgreater{}\%}
    \FunctionTok{ggplot}\NormalTok{(., }\FunctionTok{aes}\NormalTok{(}\AttributeTok{x =}\NormalTok{ Cost, }\AttributeTok{y =}\NormalTok{ TotalAmountCredited)) }\SpecialCharTok{+}
    \FunctionTok{geom\_point}\NormalTok{() }\SpecialCharTok{+}
    \FunctionTok{scale\_x\_log10}\NormalTok{() }\SpecialCharTok{+}
    \FunctionTok{labs}\NormalTok{(}\AttributeTok{x =} \StringTok{"log10(Cost)"}\NormalTok{)}

\NormalTok{x\_log\_scatter}
\end{Highlighting}
\end{Shaded}

\begin{verbatim}
## Warning in scale_x_log10(): log-10 transformation introduced infinite values.
## Removed 170 rows containing missing values or values outside the scale range
## (`geom_point()`).
\end{verbatim}

\includegraphics{lab02_files/figure-latex/unnamed-chunk-2-3.pdf}

\begin{Shaded}
\begin{Highlighting}[]
\CommentTok{\#both were pretty bad, trying them at the same time}
\NormalTok{log\_scatter }\OtherTok{\textless{}{-}}
\NormalTok{  Acres\_bmps }\SpecialCharTok{\%\textgreater{}\%}
    \FunctionTok{ggplot}\NormalTok{(., }\FunctionTok{aes}\NormalTok{(}\AttributeTok{x =}\NormalTok{ Cost, }\AttributeTok{y =}\NormalTok{ TotalAmountCredited)) }\SpecialCharTok{+}
    \FunctionTok{geom\_point}\NormalTok{() }\SpecialCharTok{+}
    \FunctionTok{scale\_x\_log10}\NormalTok{() }\SpecialCharTok{+}
    \FunctionTok{scale\_y\_log10}\NormalTok{ () }\SpecialCharTok{+}
    \FunctionTok{labs}\NormalTok{(}\AttributeTok{x =} \StringTok{"log10(Cost)"}\NormalTok{,}
         \AttributeTok{y =} \StringTok{"log10(TotalAmountCredited)"}\NormalTok{)}

\NormalTok{log\_scatter}
\end{Highlighting}
\end{Shaded}

\begin{verbatim}
## Warning in scale_x_log10(): log-10 transformation introduced infinite values.
\end{verbatim}

\begin{verbatim}
## Warning in scale_y_log10(): log-10 transformation introduced infinite values.
\end{verbatim}

\begin{verbatim}
## Warning: Removed 170 rows containing missing values or values outside the scale range
## (`geom_point()`).
\end{verbatim}

\includegraphics{lab02_files/figure-latex/unnamed-chunk-2-4.pdf}

1.3 Make a boxplot with ``StateAbbreviation'' on the x-axis and
``TotalAmountCredited'' on the y-axis. HOWEVER, the only data I want
plotted are for cover crop BMPs. Note, there are many types of cover
crops in this dataset, and I want you to include them ALL. There are
handy functions within the \texttt{stringr} package that can help you
here. Use the help function \texttt{?stringr} or read the
documentations/vignettes found online. A Google search for ``r stringr
vignettes documentation'' would be a good place to start if you're
having trouble.

\begin{Shaded}
\begin{Highlighting}[]
\CommentTok{\#filtering for only cover crop bmps}
\CommentTok{\#looking at different types of bmps}
\FunctionTok{unique}\NormalTok{(bmps}\SpecialCharTok{$}\NormalTok{BMP)}
\end{Highlighting}
\end{Shaded}

\begin{verbatim}
##   [1] "Wet Ponds and Wetlands"                                                                         
##   [2] "Dry Detention Ponds and Hydrodynamic Structures"                                                
##   [3] "Filtering Practices"                                                                            
##   [4] "Stormwater Performance Standard-Runoff Reduction"                                               
##   [5] "Bioretention/raingardens - A/B soils, underdrain"                                               
##   [6] "Bioswale"                                                                                       
##   [7] "Vegetated Open Channels - A/B soils, no underdrain"                                             
##   [8] "Permeable Pavement w/o Sand, Veg. - A/B soils, no underdrain"                                   
##   [9] "Bioretention/raingardens - C/D soils, underdrain"                                               
##  [10] "Urban Stream Restoration"                                                                       
##  [11] "Infiltration Practices w/o Sand, Veg. - A/B soils, no underdrain"                               
##  [12] "Impervious Surface Reduction"                                                                   
##  [13] "Stormwater Performance Standard-Stormwater Treatment"                                           
##  [14] "Tree Planting - Canopy"                                                                         
##  [15] "Advanced Sweeping Technology - 2 pass/week"                                                     
##  [16] "Advanced Sweeping Technology - 1 pass/week"                                                     
##  [17] "Advanced Sweeping Technology - 1 pass/2 weeks"                                                  
##  [18] "Advanced Sweeping Technology - 1 pass/4 weeks"                                                  
##  [19] "Advanced Sweeping Technology - 1 pass/8 weeks"                                                  
##  [20] "Advanced Sweeping Technology - 1 pass/12 weeks"                                                 
##  [21] "Separation of stormwater and sewer systems"                                                     
##  [22] "Erosion and Sediment Control Level 1"                                                           
##  [23] "Soil Conservation and Water Quality Plans"                                                      
##  [24] "Forest Buffer"                                                                                  
##  [25] "Tree Planting"                                                                                  
##  [26] "Land Retirement to Ag Open Space"                                                               
##  [27] "Grass Buffer"                                                                                   
##  [28] "Off Stream Watering Without Fencing"                                                            
##  [29] "Precision Intensive Rotational/Prescribed Grazing"                                              
##  [30] "Water Control Structures"                                                                       
##  [31] "Dry Extended Detention Ponds"                                                                   
##  [32] "Infiltration Practices w/ Sand, Veg. - A/B soils, no underdrain"                                
##  [33] "Forest Harvesting Practices"                                                                    
##  [34] "Cover Crop Traditional Wheat Early Drilled"                                                     
##  [35] "Cover Crop Traditional Wheat Early Other"                                                       
##  [36] "Barnyard Runoff Control"                                                                        
##  [37] "Cover Crop Traditional Wheat Normal Drilled"                                                    
##  [38] "Cover Crop Traditional Wheat Normal Other"                                                      
##  [39] "Cover Crop Traditional Wheat Late Other"                                                        
##  [40] "Cover Crop Traditional Barley Early Drilled"                                                    
##  [41] "Cover Crop Traditional Barley Early Other"                                                      
##  [42] "Cover Crop Traditional Barley Early Aerial"                                                     
##  [43] "Cover Crop Traditional Barley Normal Other"                                                     
##  [44] "Cover Crop Commodity Early"                                                                     
##  [45] "Cover Crop Commodity Normal"                                                                    
##  [46] "Cover Crop Commodity Late"                                                                      
##  [47] "Tillage Management-Conservation"                                                                
##  [48] "Nutrient Management Core N"                                                                     
##  [49] "Nutrient Management Core P"                                                                     
##  [50] "Grass Buffer-Streamside with Exclusion Fencing"                                                 
##  [51] "Cover Crop Traditional Triticale Normal Other"                                                  
##  [52] "Tillage Management-Continuous High Residue"                                                     
##  [53] "Septic Denitrification - Conventional"                                                          
##  [54] "Septic Connection"                                                                              
##  [55] "Septic Pumping"                                                                                 
##  [56] "Wetland Restoration - Floodplain"                                                               
##  [57] "Grass Buffer-Streamside with Exclusion Fencing Nitrogen Upland Acres"                           
##  [58] "Grass Buffer Nitrogen Upland Acres"                                                             
##  [59] "Forest Buffer Nitrogen Upland Acres"                                                            
##  [60] "Forest Buffer Phosphorus and Sediment Upland Acres"                                             
##  [61] "Grass Buffer Phosphorus and Sediment Upland Acres"                                              
##  [62] "Grass Buffer-Streamside with Exclusion Fencing Phosphorus and Sediment Upland Acres"            
##  [63] "Wetland Restoration Upland Acres"                                                               
##  [64] "Non Urban Stream Restoration"                                                                   
##  [65] "Cover Crop Traditional Rye Early Other"                                                         
##  [66] "Cover Crop Traditional Rye Early Aerial"                                                        
##  [67] "Cover Crop Traditional Rye Normal Drilled"                                                      
##  [68] "Cover Crop Traditional Rye Normal Other"                                                        
##  [69] "Cover Crop Traditional Rye Late Drilled"                                                        
##  [70] "Cover Crop Traditional Rye Late Other"                                                          
##  [71] "Cover Crop Traditional Wheat Early Aerial"                                                      
##  [72] "Cover Crop Traditional Wheat Late Drilled"                                                      
##  [73] "Cover Crop Traditional Barley Normal Drilled"                                                   
##  [74] "Cover Crop Traditional Forage Radish Early Other"                                               
##  [75] "Cover Crop Traditional Forage Radish Early Aerial"                                              
##  [76] "Cover Crop Traditional Annual Legume Early Other"                                               
##  [77] "Cover Crop Traditional Annual Legume Early Aerial"                                              
##  [78] "Cover Crop Traditional Triticale Early Drilled"                                                 
##  [79] "Cover Crop Traditional Triticale Early Other"                                                   
##  [80] "Cover Crop Traditional Triticale Early Aerial"                                                  
##  [81] "Cover Crop Traditional Triticale Late Drilled"                                                  
##  [82] "Cover Crop Traditional Triticale Late Other"                                                    
##  [83] "Cover Crop Traditional Annual Ryegrass Early Aerial"                                            
##  [84] "Mechanical Broom Technology - 1 pass/4 weeks"                                                   
##  [85] "Cover Crop Traditional Oats, Winter Killed Early Other"                                         
##  [86] "Cover Crop Traditional Brassica Early Other"                                                    
##  [87] "Cover Crop Traditional Brassica Early Aerial"                                                   
##  [88] "Urban Forest Buffer Upland Acres"                                                               
##  [89] "Advanced Grey Infrastructure Nutrient Discovery Program (IDDE)"                                 
##  [90] "Grass Buffer-Narrow with Exclusion Fencing"                                                     
##  [91] "Land Retirement to Pasture"                                                                     
##  [92] "Forest Buffer-Streamside with Exclusion Fencing"                                                
##  [93] "Horse Pasture Management"                                                                       
##  [94] "Impervious Disconnection to amended soils"                                                      
##  [95] "Cover Crop Traditional Rye Early Drilled"                                                       
##  [96] "Forest Planting"                                                                                
##  [97] "Loafing Lot Management"                                                                         
##  [98] "Manure Injection"                                                                               
##  [99] "Manure Incorporation Low Disturbance Late"                                                      
## [100] "Nutrient Management Maryland Commercial Applicators"                                            
## [101] "Nutrient Management Maryland Do It Yourself"                                                    
## [102] "Cover Crop Traditional Triticale Normal Drilled"                                                
## [103] "Erosion and Sediment Control Level 2"                                                           
## [104] "Forest Buffer-Streamside with Exclusion Fencing Nitrogen Upland Acres"                          
## [105] "Forest Buffer-Streamside with Exclusion Fencing Phosphorus and Sediment Upland Acres"           
## [106] "Cover Crop Traditional Oats, Winter Hardy Early Drilled"                                        
## [107] "Cover Crop Traditional Oats, Winter Hardy Early Aerial"                                         
## [108] "Cover Crop Traditional Brassica Early Drilled"                                                  
## [109] "Wetland Creation - Floodplain"                                                                  
## [110] "Wetland Creation Upland Acres"                                                                  
## [111] "Forest Buffer-Narrow with Exclusion Fencing"                                                    
## [112] "Urban Shoreline Management"                                                                     
## [113] "Dirt & Gravel Road Erosion & Sediment Control - Driving Surface Aggregate + Raising the Roadbed"
## [114] "Non Urban Shoreline Management"                                                                 
## [115] "Cover Crop Traditional Annual Ryegrass Early Drilled"                                           
## [116] "Cover Crop Traditional Oats, Winter Hardy Normal Drilled"                                       
## [117] "Alternative Crops"                                                                              
## [118] "Abandoned Mine Reclamation"                                                                     
## [119] "Cover Crop Traditional with Fall Nutrients Rye Normal Drilled"                                  
## [120] "Tillage Management-Low Residue"                                                                 
## [121] "Nutrient Management N Rate"                                                                     
## [122] "Nutrient Management P Rate"                                                                     
## [123] "Nutrient Management N Placement"                                                                
## [124] "Nutrient Management P Placement"                                                                
## [125] "Nutrient Management N Timing"                                                                   
## [126] "Nutrient Management P Timing"                                                                   
## [127] "Forest Buffer - Narrow"                                                                         
## [128] "Grass Buffer - Narrow"                                                                          
## [129] "Cover Crop Traditional with Fall Nutrients Rye Normal Other"                                    
## [130] "Wetland Enhancement"                                                                            
## [131] "Wetland Enhancement Upland Acres"                                                               
## [132] "Non Urban Shoreline Erosion Control Non-Vegetated"                                              
## [133] "Urban Shoreline Erosion Control Non-Vegetated"                                                  
## [134] "Cover Crop Traditional Annual Legume Normal Drilled"                                            
## [135] "Nutrient Management Plan"                                                                       
## [136] "Permeable Pavement w/o Sand, Veg. - C/D soils, underdrain"                                      
## [137] "Septic Secondary Treatment - Enhanced"                                                          
## [138] "Septic Secondary Treatment - Conventional"                                                      
## [139] "Vegetated Open Channels - C/D soils, no underdrain"                                             
## [140] "Permeable Pavement w/ Sand, Veg. - A/B soils, no underdrain"                                    
## [141] "Permeable Pavement w/ Sand, Veg. - C/D soils, underdrain"                                       
## [142] "Cover Crop Traditional with Fall Nutrients Wheat Early Other"                                   
## [143] "Septic Denitrification - Enhanced"                                                              
## [144] "Bioretention/raingardens - A/B soils, no underdrain"                                            
## [145] "Septic Effluent - Enhanced"                                                                     
## [146] "Urban Shoreline Erosion Control Vegetated"                                                      
## [147] "Permeable Pavement w/ Sand, Veg. - A/B soils, underdrain"                                       
## [148] "Permeable Pavement w/o Sand, Veg. - A/B soils, underdrain"                                      
## [149] "Filter Strip Runoff Reduction"                                                                  
## [150] "Filter Strip Stormwater Treatment"                                                              
## [151] "Cover Crop Traditional with Fall Nutrients Rye Early Other"                                     
## [152] "Cover Crop Traditional with Fall Nutrients Wheat Normal Other"                                  
## [153] "Cover Crop Traditional Forage Radish Early Drilled"                                             
## [154] "Cover Crop Traditional Legume Plus Grass 25-50% Early Drilled"                                  
## [155] "Cover Crop Traditional Legume Plus Grass 25-50% Early Other"                                    
## [156] "Cover Crop Traditional Annual Legume Early Drilled"                                             
## [157] "Cover Crop Traditional Forage Radish Plus Early Drilled"                                        
## [158] "Cover Crop Traditional Forage Radish Plus Normal Drilled"                                       
## [159] "Cover Crop Traditional Legume Plus Grass 25-50% Normal Drilled"
\end{verbatim}

\begin{Shaded}
\begin{Highlighting}[]
\NormalTok{cover\_crop }\OtherTok{\textless{}{-}}
\NormalTok{  bmps }\SpecialCharTok{|\textgreater{}} 
  \FunctionTok{filter}\NormalTok{(}\FunctionTok{str\_detect}\NormalTok{(BMP, }\StringTok{"Cover Crop"}\NormalTok{)) }\SpecialCharTok{|\textgreater{}}
  \FunctionTok{ggplot}\NormalTok{(}\FunctionTok{aes}\NormalTok{(}\AttributeTok{x =}\NormalTok{ StateAbbreviation, }\AttributeTok{y =}\NormalTok{ TotalAmountCredited)) }\SpecialCharTok{+}
    \FunctionTok{geom\_boxplot}\NormalTok{(}\FunctionTok{aes}\NormalTok{(}\AttributeTok{fill =}\NormalTok{ StateAbbreviation)) }

\NormalTok{cover\_crop}
\end{Highlighting}
\end{Shaded}

\includegraphics{lab02_files/figure-latex/unnamed-chunk-3-1.pdf}

1.4 make a scatterplot of the \texttt{dam} dataset, this time with
``YEAR'' on the x-axis and ``STATE'' on y-axis (think of it like a
timeline). Assume no dams were built in year 0, so you'll need to remove
those data points. You know how to remove observations based on a given
criteria.

\begin{Shaded}
\begin{Highlighting}[]
\CommentTok{\#remove year 0 from dams and plot!}
\NormalTok{dams\_new }\OtherTok{\textless{}{-}}\NormalTok{ dams }\SpecialCharTok{|\textgreater{}}
  \FunctionTok{filter}\NormalTok{(YEAR }\SpecialCharTok{!=} \DecValTok{0}\NormalTok{) }\SpecialCharTok{|\textgreater{}}
  \FunctionTok{ggplot}\NormalTok{(}\FunctionTok{aes}\NormalTok{(}\AttributeTok{x =}\NormalTok{ YEAR, }\AttributeTok{y =}\NormalTok{ STATE)) }\SpecialCharTok{+}
    \FunctionTok{geom\_point}\NormalTok{() }
  
\NormalTok{dams\_new}
\end{Highlighting}
\end{Shaded}

\includegraphics{lab02_files/figure-latex/unnamed-chunk-4-1.pdf}

1.5 make one last (aspatial) visualization. But this time, it's your
choice what data and plots to use. The only requirement is that you link
two of the datasets together in some manner. Be creative. Make it look
nice (e.g., use proper labels, interesting colors/shading/size).

\subsubsection{Task 2: Spatial
operations}\label{task-2-spatial-operations}

2.1 Find the 5 longest streams in the `streams opened by dam removal'
dataset

\begin{Shaded}
\begin{Highlighting}[]
\CommentTok{\#group by GNIS name}
\NormalTok{stream\_length}\OtherTok{\textless{}{-}}\NormalTok{ streams }\SpecialCharTok{|\textgreater{}} \FunctionTok{group\_by}\NormalTok{(GNIS\_Name) }\SpecialCharTok{|\textgreater{}} 
  \FunctionTok{filter}\NormalTok{(}\SpecialCharTok{!}\FunctionTok{is.na}\NormalTok{(GNIS\_Name)) }\SpecialCharTok{|\textgreater{}} \CommentTok{\#remove unnamed rivers}
  \FunctionTok{summarize}\NormalTok{(}\AttributeTok{TotalLength=}\FunctionTok{sum}\NormalTok{(LengthKM)) }
  

\NormalTok{stream\_length}
\end{Highlighting}
\end{Shaded}

\begin{verbatim}
## Simple feature collection with 113 features and 2 fields
## Geometry type: GEOMETRY
## Dimension:     XY
## Bounding box:  xmin: -79.21001 ymin: 37.23123 xmax: -75.53429 ymax: 41.75646
## Geodetic CRS:  WGS 84
## # A tibble: 113 x 3
##    GNIS_Name        TotalLength                                         geometry
##    <chr>                  <dbl>                                   <GEOMETRY [°]>
##  1 Adams Run              1.16  LINESTRING (-76.39092 40.59092, -76.39099 40.59~
##  2 Appomattox River       0.466 LINESTRING (-77.4212 37.23123, -77.42115 37.231~
##  3 Back Creek            14.9   LINESTRING (-78.98025 37.92897, -78.97991 37.92~
##  4 Barr Run               2.15  LINESTRING (-78.266 41.56257, -78.26632 41.5624~
##  5 Barren Creek          11.2   LINESTRING (-75.63235 38.4521, -75.63236 38.452~
##  6 Big Run                5.48  MULTILINESTRING ((-78.36152 41.56917, -78.36288~
##  7 Billy Buck Run         2.44  LINESTRING (-78.42342 41.59632, -78.42362 41.59~
##  8 Blackwater Creek      20.0   MULTILINESTRING ((-75.62309 38.4635, -75.6231 3~
##  9 Bobby Run              8.10  LINESTRING (-78.41307 41.5404, -78.41263 41.540~
## 10 Brewer Fork            3.13  LINESTRING (-78.37411 41.6047, -78.37444 41.604~
## # i 103 more rows
\end{verbatim}

2.2 Find the three counties with the greatest TOTAL length of streams
(opened by dam removal) in them

\begin{Shaded}
\begin{Highlighting}[]
\CommentTok{\#find what counties contain the highest number of streams}
\NormalTok{counties\_streams }\OtherTok{\textless{}{-}} \FunctionTok{st\_join}\NormalTok{(counties, streams, }\AttributeTok{join =}\NormalTok{ st\_intersects, }\AttributeTok{left=}\ConstantTok{FALSE}\NormalTok{) }\CommentTok{\#inner join}

\CommentTok{\# all counties and rivers}
\FunctionTok{tm\_shape}\NormalTok{(counties) }\SpecialCharTok{+} \FunctionTok{tm\_polygons}\NormalTok{(}\AttributeTok{fill=}\StringTok{"grey"}\NormalTok{) }\SpecialCharTok{+} \FunctionTok{tm\_shape}\NormalTok{(streams) }\SpecialCharTok{+} \FunctionTok{tm\_lines}\NormalTok{(}\AttributeTok{col =} \StringTok{"blue"}\NormalTok{)}
\end{Highlighting}
\end{Shaded}

\includegraphics{lab02_files/figure-latex/unnamed-chunk-7-1.pdf}

\begin{Shaded}
\begin{Highlighting}[]
\CommentTok{\#quick map to check}
\FunctionTok{tm\_shape}\NormalTok{(counties) }\SpecialCharTok{+} \FunctionTok{tm\_polygons}\NormalTok{(}\AttributeTok{fill=}\StringTok{"grey"}\NormalTok{) }\SpecialCharTok{+} \FunctionTok{tm\_shape}\NormalTok{(counties\_streams) }\SpecialCharTok{+} \FunctionTok{tm\_lines}\NormalTok{(}\AttributeTok{col =} \StringTok{"blue"}\NormalTok{)}
\end{Highlighting}
\end{Shaded}

\includegraphics{lab02_files/figure-latex/unnamed-chunk-7-2.pdf}

\begin{Shaded}
\begin{Highlighting}[]
\CommentTok{\#top 3 counties with greatest total length of streams}
\NormalTok{counties\_streams }\SpecialCharTok{|\textgreater{}} \FunctionTok{group\_by}\NormalTok{(COUNTYFP10) }\SpecialCharTok{|\textgreater{}}
  \FunctionTok{summarize}\NormalTok{(}\AttributeTok{AllStreamLength=}\FunctionTok{sum}\NormalTok{(LengthKM)) }\SpecialCharTok{|\textgreater{}}
  \FunctionTok{slice\_max}\NormalTok{(AllStreamLength, }\AttributeTok{n=}\DecValTok{3}\NormalTok{)}
\end{Highlighting}
\end{Shaded}

\begin{verbatim}
## Simple feature collection with 3 features and 2 fields
## Geometry type: POLYGON
## Dimension:     XY
## Bounding box:  xmin: -79.53331 ymin: 37.88159 xmax: -77.67661 ymax: 41.61821
## Geodetic CRS:  WGS 84
## # A tibble: 3 x 3
##   COUNTYFP10 AllStreamLength                                            geometry
##   <chr>                <dbl>                                       <POLYGON [°]>
## 1 015                   425. ((-79.36744 38.0179, -79.35294 38.00927, -79.34872~
## 2 061                   309. ((-78.14963 40.1743, -78.14926 40.17408, -78.14071~
## 3 023                   242. ((-78.25396 41.40335, -78.25398 41.39962, -78.2540~
\end{verbatim}

2.3 Make a map of the counties, shading each county by the total cost of
BMPs funded/implemented in that county. This will require you to join
multiple datasets together

2.4 For each removed dam, find the closest stream segment

2.5 Calculate how many removed dams are (or were) in each state

\end{document}
